% !PRINTED 630
\weberpage{630}
%<<CONT
prime factors conjugate to $\mathfrak p_1$, and it follows,
since $p$ is not divisible by $\mathfrak p^2$:

1. \gesperrt{The prime factor $\mathfrak p_1$ is the greatest common
divisor of $p$ and $P_1(r)$; in the same way it follows that
$\mathfrak p_n$ is the greatest common divisor of $p$ and $P_1(r^n)$,
or also the greatest common divisor of $p$ and
$P_{n'}(r)$, when $n n' \equiv 1 \ (\mathrm{mod}\ m)$.}

The latter follows from the congruence which results from (4)
by the substitution $(r, r^{n'})$,
\[
F_1(t) \equiv P_n(t) \ (\mathrm{mod}\ \mathfrak p_{n'}).
\]

One thus obtains all the prime factors of $p$ if one
seeks out the greatest common divisors of $p$ with each
of the rational functions
\weq{10}{P_{\xi_1}(r),\ P_{\xi_2}(r),\ \ldots,\ P_{\xi_e}(r)}
and indeed $\mathfrak p_{\xi_i}$ is the greatest common divisor
of $p$ and $P_{\xi_i^{-1}}$.
% [sic] Source omits the argument (r) after P_{xi_i^{-1}} here; (10) and theorem 1 give the intended evaluation.

We shall consider somewhat more closely the case where $f=1$,
and thus
\weq{11}{p \equiv 1 \ (\mathrm{mod}\ m).}

In this case $e = \varphi(m)$; the group $\mathfrak A$ reduces
to the unit, and all the $\mu$ conjugate factors $\mathfrak p$
of $p$ are different from one another. The forms
\[
P_{\xi_1}(t),\ P_{\xi_2}(t),\ \ldots,\ P_{\xi_e}(t)
\]
all become linear, that is, $r$ is congruent to a rational number
with respect to each of the moduli $\mathfrak p_i$.

This follows also from the fact that here the number of
numbers incongruent with respect to the modulus $\mathfrak p$ is equal to $N(\mathfrak p) = p$,
and that therefore $0, 1, 2, \ldots, p-1$ is a complete system of residues.

If accordingly, say, $r \equiv c \ (\mathrm{mod}\ \mathfrak p)$, then, since $r^m = 1$,
we must have $c^m \equiv 1 \ (\mathrm{mod}\ \mathfrak p)$ and hence also $(\mathrm{mod}\ p)$, and if
now $g$ is a primitive root of the prime $p$, then we must have
\weq{12}{c \equiv g^{-n\frac{p-1}{m}} \ (\mathrm{mod}\ p)}
in which $n$ is relatively prime to $m$. If we let $\mathfrak p$ run through
the system of conjugate ideals, then $n$ in (12) must run through the
group $\mathfrak N$, and it is therefore only a matter of notation
if we stipulate:
\weq{13}{r \equiv g^{-\frac{p-1}{m}} \ (\mathrm{mod}\ \mathfrak p).}
